\section[Carnet 1968-14 : préimages compactes des fermés loin de $A$]{Carnet 1968-14 : applications dont les préimages des fermés loin de $A$ sont compactes\protect\piste\protect\verifie{Notebook1968_14}}\label{sec:quillen-1968-14}

\emph{Le contexte.} Le \carnet{1968}{14}, de décembre 1968, cherche une description géométrique de l'homologie, par classes de bordisme d'applications de variétés orientées. Pour l'homologie d'un espace $X$ relative à un fermé $A$, les cycles proposés aux pages~26--27 sont des applications $f : V \to X \setminus A$ telles que $f^{-1}(B)$ soit compact pour tout fermé $B$ de $X$ disjoint de $A$. Quillen note que, si $X$ est compact, cette condition revient à dire que $f$ est propre comme application vers $X \setminus A$ : un fermé de $X$ disjoint de $A$ est alors un compact de $X \setminus A$, et réciproquement.

\emph{La page.} À la page~27, il énonce une proposition sur les limites inductives prises sur les compacts de $A$. Le cœur de la preuve est le point suivant : si $f$ vérifie la condition, l'intersection de l'adhérence de $f(V)$ avec $A$ est compacte. Il raisonne par l'absurde : un fermé non compact contient une suite discrète $(a_n)$ ; on choisit une suite exhaustive $(Q_n)$ de compacts de $V$, puis des points $x_n \notin Q_n$ dont les images $f(x_n)$ sont assez proches des $a_n$ pour former un ensemble discret. Cet ensemble est fermé, disjoint de $A$, et sa préimage n'est contenue dans aucun $Q_n$ : elle n'est pas compacte. En haut de la page~28, une remarque tire de la preuve l'équivalence de la condition avec deux propriétés : $f$ est propre, et l'adhérence de $f(V)$ dans $X$ est compacte. C'est cette remarque que l'on démontre ici. Ses espaces sont des variétés ; la preuve de la page se sert d'une exhaustion de $V$ par des compacts, donc de la $\sigma$-compacité de $V$.

\emph{L'intuition.} Une application propre vers $X \setminus A$ peut encore \q{partir à l'infini} de deux façons, vues depuis $X$ : en s'approchant de $A$, ce que la condition autorise, ou en sortant de tout compact de $X$, ce qu'elle interdit. La condition dit exactement que $f$ ne fuit que vers $A$ (figure~\ref{fig:quillen-1968-14}). La compacité de l'adhérence de $f(V)$ exclut la seconde fuite, la propreté vers $X \setminus A$ dit que la première est la seule.


\begin{figure}[htbp]
\centering
\begin{tikzpicture}[every node/.style={font=\small}]
\draw[black!55, rounded corners=12pt] (0,0) rectangle (9,4.4);
\node[black!65] at (8.6,4.0) {$X$};
\fill[frise!14] (2.901,1.340) .. controls (4,0.9) and (6,0.8) .. (7,1.6) .. controls (7.6,2.4) and (6.5,3.3) .. (5,3.1) .. controls (4,3.0) and (3.4,3.2) .. (2.901,3.060) arc[start angle=35, end angle=-35, x radius=1.1, y radius=1.5];
\draw[frise] (2.901,1.340) .. controls (4,0.9) and (6,0.8) .. (7,1.6) .. controls (7.6,2.4) and (6.5,3.3) .. (5,3.1) .. controls (4,3.0) and (3.4,3.2) .. (2.901,3.060);
\fill[piste!20] (2,2.2) ellipse [x radius=1.1, y radius=1.5];
\draw[piste] (2,2.2) ellipse [x radius=1.1, y radius=1.5];
\draw[piste, line width=1.6pt] (2.901,3.060) arc[start angle=35, end angle=-35, x radius=1.1, y radius=1.5];
\node[piste] at (1.7,2.2) {$A$};
\node[frise] at (4.6,2.35) {$f(V)$};
\fill[black!35, opacity=0.35] (6.2,2.0) ellipse [x radius=0.8, y radius=0.6];
\draw[black!70, dashed] (6.2,2.0) ellipse [x radius=0.8, y radius=0.6];
\node at (6.35,2.0) {$B$};
\end{tikzpicture}
\caption{Les conditions du théorème~\ref{thm:quillen-1968-14}. L'image $f(V)$ (en bleu) évite $A$ (en rouge) mais s'en approche : son adhérence rencontre $A$, ici le long de l'arc en trait épais. La condition~(1) demande que la préimage de tout fermé $B$ disjoint de $A$ (en gris) soit compacte : $f$ peut fuir vers $A$, et seulement vers $A$. La condition~(2) dit la même chose autrement : $f$ est propre vers $X \setminus A$, et l'adhérence de $f(V)$ est compacte.}
\label{fig:quillen-1968-14}
\end{figure}

\begin{encadre}
\begin{theoreme}[d'après Quillen, carnet 1968-14, p.~28]\label{thm:quillen-1968-14}
Soient $X$ un espace métrisable, $A$ une partie de $X$, $V$ un espace topologique et $f : V \to X$ une application continue à valeurs dans $X \setminus A$. Les conditions suivantes sont équivalentes :
\begin{enumerate}
\item pour tout fermé $B$ de $X$ disjoint de $A$, la partie $f^{-1}(B)$ de $V$ est compacte ;
\item $f$ est propre comme application de $V$ dans $X \setminus A$, et l'adhérence de $f(V)$ dans $X$ est compacte.
\end{enumerate}
\end{theoreme}
\end{encadre}

\noindent Les hypothèses sont plus faibles que celles de la page : $X$ n'est pas supposé localement compact, $A$ n'est pas supposé fermé, et $V$ est quelconque, sans $\sigma$-compacité. La preuve ci-dessous remplace l'exhaustion de $V$ par des suites de valeurs de $f$. Seule la métrisabilité de $X$ sert, dans le sens (1) $\Rightarrow$ (2), et seulement pour pouvoir raisonner sur des suites.

\begin{proof}
\emph{(2) $\Rightarrow$ (1).} Soit $K$ l'adhérence de $f(V)$, compacte. Pour un fermé $B$ disjoint de $A$, la partie $B \cap K$ est compacte, comme fermé d'un compact, et contenue dans $X \setminus A$ ; c'est donc un compact de $X \setminus A$. Comme $f(V) \subset K$, on a $f^{-1}(B) = f^{-1}(B \cap K)$, préimage d'un compact par une application propre : elle est compacte. Ce sens n'utilise aucune hypothèse sur $X$.

\emph{Un lemme sur les suites.} Supposons (1). Soit $(v_n)$ une suite de $V$ ; montrons que la suite $(f(v_n))$ a une valeur d'adhérence dans $X$. Sinon, soit $B$ l'ensemble de ses valeurs. Il est fermé : un point $y$ de l'adhérence de $B$ qui n'est pas dans $B$ serait valeur d'adhérence de la suite, car tout voisinage de $y$, privé des valeurs en nombre fini $f(v_0), \dots, f(v_{N-1})$ (qui forment un fermé ne contenant pas $y$), rencontre encore $B$, donc contient un $f(v_n)$ avec $n \ge N$. Il est disjoint de $A$, puisque $f$ évite $A$. Par (1), $f^{-1}(B)$ est compact ; il contient tous les $v_n$, qui ont donc une valeur d'adhérence $w$. Par continuité, $f(w)$ est valeur d'adhérence de $(f(v_n))$ : contradiction.

\emph{(1) $\Rightarrow$ adhérence compacte.} On munit $X$ d'une distance. Dans un espace métrique, une partie est compacte si et seulement si toute suite de cette partie a une sous-suite convergente dans la partie. Soit $(z_n)$ une suite de l'adhérence $K$ de $f(V)$. On choisit $v_n \in V$ avec $d(z_n, f(v_n)) < 1/(n+1)$. Par le lemme, $(f(v_n))$ a une valeur d'adhérence $y$, donc une sous-suite $(f(v_{\varphi(k)}))$ qui tend vers $y$ (un espace métrique est à base dénombrable de voisinages). Comme $\varphi(k) \ge k$, la distance de $z_{\varphi(k)}$ à $f(v_{\varphi(k)})$ est majorée par $1/(k+1)$, et $z_{\varphi(k)}$ tend aussi vers $y$. Enfin $y$ est limite de points de $f(V)$, donc dans $K$.

\emph{(1) $\Rightarrow$ propreté.} L'espace $X \setminus A$ est métrisable, donc séparé et à génération compacte ; dans un tel espace d'arrivée, une application continue est propre si et seulement si la préimage de tout compact est compacte. Soit $C$ un compact de $X \setminus A$. Vu dans $X$, c'est un compact, donc un fermé ($X$ est séparé), et il est disjoint de $A$. Sa préimage est compacte par (1).
\end{proof}

\noindent La preuve de la page passe par l'intersection de l'adhérence de $f(V)$ avec $A$ ; celle-ci montre directement que l'adhérence entière est compacte, ce que la remarque de la page affirme. Le choix d'une suite discrète, que Quillen fait dans $A$ et approche par des valeurs de $f$, est fait ici directement parmi les valeurs de $f$, ce qui dispense de l'exhaustion de $V$.

\begin{exemple}
Soient $X = \mathbf R$, $A = \{0\}$ et $f$ l'inclusion de $V = \,]0, 1[$. Le fermé $B = [1/2, 1]$ ne rencontre pas $A$, et $f^{-1}(B) = [1/2, 1[$ n'est pas compact : la condition (1) échoue, et $f$ n'est pas propre vers $\mathbf R \setminus \{0\}$, bien que l'adhérence $[0, 1]$ de $f(V)$ soit compacte. Avec $V = \,]0, 1]$, au contraire, $f^{-1}(B) = B \cap [0, 1]$ est compact pour tout fermé $B$ évitant $0$ : les deux conditions sont vraies, l'application fuit vers $A = \{0\}$ et seulement vers lui. Avec $A = \emptyset$, la condition (1) appliquée à $B = X$ dit que $V$ est compact, et (2) le dit aussi.
\end{exemple}

\emph{Vérification en Lean.} Le fichier \href{\leanblobmain/Notebook1968_14.lean}{\texttt{Notebook1968\_14.\allowbreak{}lean}} démontre le théorème~\ref{thm:quillen-1968-14} sous le nom \texttt{preimage\_compact\_iff}, pour $X$ muni de la classe \texttt{MetrizableSpace} et $f$ restreinte au but $A^c$ par \texttt{Set.codRestrict}. Le lemme sur les suites est \texttt{exists\_mapClusterPt}, appuyé sur \texttt{mapClusterPt\_of\_mem\_closure\_range} ; la compacité de l'adhérence est \texttt{isCompact\_closure\_range}, la propreté \texttt{isProperMap\_codRestrict}, par le critère de mathlib \texttt{isProperMap\_iff\_isCompact\_preimage} ; la réciproque est \texttt{isCompact\_preimage\_of\_isProperMap}, énoncée pour un espace $X$ quelconque. Les axiomes utilisés sont les trois axiomes usuels de mathlib (\texttt{propext}, \texttt{Classical.choice}, \texttt{Quot.sound}).

Ne sont pas démontrés ici : la proposition de la page~27 sur les limites inductives d'homologie, dont l'énoncé est le point de départ ; la description de l'homologie par bordisme ; la question de savoir si la métrisabilité de $X$ peut être remplacée par une hypothèse plus faible (espace séquentiel, par exemple).

\litterature{la recherche n'a pas retrouvé cet énoncé imprimé sous cette forme. Il est élémentaire et voisin de résultats classiques sur les applications propres (Bourbaki, \emph{Topologie générale}, chapitre~I, \S\,10) ; le cadre, homologie relative par bordisme, se rattache aux travaux de Quillen sur le cobordisme publiés en 1969--1971. L'édition ne le donne pas pour nouveau. Ce n'est pas une piste.}
